\section{Implementacja}		%4

\hspace{0.60cm}Niniejszy rozdział przedstawia szczegóły implementacyjne aplikacji mobilnej „Guess the Beat”. Omówiono w nim strukturę pakietu \texttt{com.example.guessthebeat}, zastosowane dziedziczenie, autentyczne fragmenty kodu źródłowego wszystkich kluczowych klas, repozytoriów chmurowych, modułów sprzętowych i widoków oraz przebieg zmian wdrożonych w trakcie prac projektowych.

\subsection{Struktura plików źródłowych i ich odpowiedzialność}

\hspace{0.60cm}Kod źródłowy projektu został podzielony na wyspecjalizowane pakiety i pliki zgodne z architekturą MVVM (Model-View-ViewModel):
\begin{itemize}
    \item \textbf{\texttt{MainActivity.kt}} – punkt wejścia aplikacji. Dziedziczy po \texttt{FragmentActivity}, konfiguruje ciemny motyw z neonowym gradientem (\texttt{NeonBlue}), zarządza cyklem życia odtwarzacza audio oraz spina nawigację między ekranami za pomocą komponentu \texttt{Scaffold}.
    \item \textbf{\texttt{viewmodel/GameViewModel.kt}} – główny kontroler stanu aplikacji. Przechowuje aktualny wynik punktowy, numer rundy, listy przebojów (\textit{Charts}), historię błędnych odpowiedzi, zarządza oknem potwierdzenia wyjścia z gry oraz komunikacją z warstwą sieciową.
    \item \textbf{\texttt{data/firebase/FirebaseRepositories.kt}} – warstwa integracji z usługami chmurowymi Google. Zawiera model statystyk \texttt{PlayerStats}, klasę \texttt{AuthRepository} (obsługującą logowanie e-mailem, rejestrację oraz logowanie przez konto Google) oraz klasę \texttt{StatsRepository} (zapisującą i odczytującą wyniki gracza w bazie Cloud Firestore).
    \item \textbf{\texttt{data/remote/DeezerAPI.kt}} – warstwa sieciowa i multimedialna. Zawiera modele danych (\texttt{SongSuggestion}, \texttt{ChartItem}) oraz klasę \texttt{MusicPlayer}, odpowiedzialną za zapytania HTTP do API Deezer i sterowanie obiektem \texttt{MediaPlayer}.
    \item \textbf{\texttt{sensors/BiometricAuthManager.kt}} – moduł bezpieczeństwa zarządzający dostępnością czujników biometrycznych urządzenia, bezpiecznym przechowywaniem poświadczeń w \texttt{SharedPreferences} oraz wyświetlaniem systemowego okna \texttt{BiometricPrompt}.
    \item \textbf{\texttt{sensors/FlipDetector.kt}} – moduł sprzętowy implementujący interfejs \texttt{SensorEventListener}, odpowiedzialny za detekcję obrotu smartfona ekranem do dołu za pomocą akcelerometru.
    \item \textbf{\texttt{ui/screens/MainMenuScreen.kt}} – deklaratywny widok menu głównego, zawierający wyszukiwarkę playlist, polecane kategorie muzyczne oraz karuzele aktualnych list przebojów (\textit{Top Playlisty, Artyści, Albumy}).
    \item \textbf{\texttt{ui/screens/GameScreen.kt}} – deklaratywny widok rozgrywki, obsługujący stany rundy, pasek postępu prób, listę podpowiedzi, ekran podsumowania sesji oraz dialog potwierdzenia opuszczenia trwającej gry.
    \item \textbf{\texttt{ui/screens/LoginScreen.kt}} – deklaratywny ekran uwierzytelniania, umożliwiający logowanie, rejestrację konta, logowanie przez Google oraz wywołanie autoryzacji biometrycznej.
    \item \textbf{\texttt{ui/components/BottomNavBar.kt}} – dolny pasek nawigacyjny umożliwiający przełączanie się między ekranem menu głównego a panelem konta gracza.
\end{itemize}

\subsection{Dziedziczenie i reaktywne zarządzanie stanem}

\hspace{0.60cm}Architektura systemu opiera się na mechanizmach dziedziczenia dostarczanych przez platformę Android:
\begin{itemize}
    \item Klasa \texttt{MainActivity} dziedziczy po \texttt{FragmentActivity}, co zapewnia pełną kompatybilność zarówno z drzewem widoków Jetpack Compose (\texttt{setContent}), jak i z systemowym kontrolerem okna biometrycznego \texttt{BiometricPrompt}.
    \item Klasa \texttt{GameViewModel} dziedziczy po \texttt{androidx.lifecycle.ViewModel}. Wykorzystuje delegaty \texttt{mutableStateOf} oraz \texttt{mutableStateListOf}, dzięki czemu każda zmiana wartości zmiennej (np. aktualizacja punktacji, zmiana aktywnego ekranu czy wyświetlenie dialogu wyjścia) automatycznie wyzwala rekompozycję odpowiednich elementów interfejsu.
    \item Klasa \texttt{FlipDetector} implementuje interfejs \texttt{SensorEventListener}, umożliwiając nasłuchiwanie zdarzeń sprzętowych czujnika ruchu.
\end{itemize}

\subsection{Główna aktywność i nawigacja reaktywna (MainActivity.kt)}

\hspace{0.60cm}W pliku \texttt{MainActivity.kt} zaimplementowano główny kontener widoku (\texttt{MainContainer}) oparty na strukturze \texttt{Scaffold} z dolnym paskiem nawigacyjnym \texttt{BottomNavBar}. Zamiast tradycyjnego przełączania wielu aktywności za pomocą obiektów \texttt{Intent}, aplikacja dynamicznie podmienia funkcje kompozywalne na podstawie wartości zmiennych \texttt{viewModel.currentScreen} oraz \texttt{viewModel.selectedPlaylist}, integrując jednocześnie obsługę logowania Firebase i biometrii.

\begin{lstlisting}[language=Kotlin, caption={Implementacja klasy MainActivity oraz głównego kontenera nawigacji w pliku MainActivity.kt}, label={lst:main_activity}]
class MainActivity : FragmentActivity() {
    private val musicPlayer = MusicPlayer()

    override fun onCreate(savedInstanceState: Bundle?) {
        super.onCreate(savedInstanceState)
        setContent {
            GuessTheBeatTheme {
                val topColor = Color(0xFF070707)
                val bottomNeon = NeonBlue.copy(alpha = 0.62f)

                val backgroundGradient = Brush.verticalGradient(
                    colors = listOf(topColor, bottomNeon)
                )

                Surface(
                    modifier = Modifier
                        .fillMaxSize()
                        .background(backgroundGradient),
                    color = Color.Transparent
                ) {
                    MainContainer(musicPlayer = musicPlayer)
                }
            }
        }
    }

    override fun onPause() {
        super.onPause()
        musicPlayer.stop()
    }
}

@Composable
fun MainContainer(
    musicPlayer: MusicPlayer,
    viewModel: GameViewModel = viewModel()
) {
    val context = LocalContext.current
    val scope = rememberCoroutineScope()
    val authRepository = remember { AuthRepository() }
    val statsRepository = remember { StatsRepository() }
    val biometricAuthManager = remember(context) { BiometricAuthManager(context) }

    var isLoading by remember { mutableStateOf(false) }
    var errorMessage by remember { mutableStateOf<String?>(null) }
    var currentUserEmail by remember { mutableStateOf(authRepository.currentUserEmail) }
    var playerStats by remember { mutableStateOf(PlayerStats()) }

    LaunchedEffect(authRepository.currentUserId) {
        val uid = authRepository.currentUserId
        if (uid != null) {
            runCatching { statsRepository.getStats(uid) }
                .onSuccess { playerStats = it }
        } else {
            playerStats = PlayerStats()
        }
    }

    Scaffold(
        bottomBar = {
            BottomNavBar(viewModel = viewModel, musicPlayer = musicPlayer)
        },
        containerColor = Color.Transparent
    ) { innerPadding ->
        Column(
            modifier = Modifier
                .fillMaxSize()
                .padding(top = innerPadding.calculateTopPadding())
                .padding(horizontal = 16.dp),
            horizontalAlignment = Alignment.CenterHorizontally
        ) {
            Spacer(modifier = Modifier.height(16.dp))
            Image(
                painter = painterResource(id = R.drawable.logo),
                contentDescription = "Logo",
                modifier = Modifier.size(120.dp)
            )
            Text(
                text = "GuessTheBeat",
                style = MaterialTheme.typography.titleLarge.copy(
                    fontWeight = FontWeight.Bold,
                    color = MaterialTheme.colorScheme.primary
                ),
                modifier = Modifier.padding(top = 4.dp, bottom = 2.dp)
            )

            when {
                viewModel.currentScreen == "login" -> {
                    LoginScreen(
                        isLoading = isLoading,
                        errorMessage = errorMessage,
                        currentUserEmail = currentUserEmail,
                        isBiometricAvailable = biometricAuthManager.isBiometricAvailable(),
                        onLogin = { email, password ->
                            scope.launch {
                                isLoading = true
                                errorMessage = null
                                authRepository.login(email, password)
                                    .onSuccess {
                                        biometricAuthManager.saveCredentials(email, password)
                                        currentUserEmail = authRepository.currentUserEmail
                                        viewModel.currentScreen = "menu"
                                    }
                                    .onFailure { e ->
                                        errorMessage = e.localizedMessage ?: "Błąd logowania"
                                    }
                                isLoading = false
                            }
                        },
                        onRegister = { email, password ->
                            scope.launch {
                                isLoading = true
                                errorMessage = null
                                authRepository.register(email, password)
                                    .onSuccess {
                                        biometricAuthManager.saveCredentials(email, password)
                                        currentUserEmail = authRepository.currentUserEmail
                                        viewModel.currentScreen = "menu"
                                    }
                                    .onFailure { e ->
                                        errorMessage = e.localizedMessage ?: "Błąd rejestracji"
                                    }
                                isLoading = false
                            }
                        },
                        onGoogleLogin = {
                            scope.launch {
                                isLoading = true
                                errorMessage = null
                                authRepository.loginWithGoogle(context)
                                    .onSuccess {
                                        currentUserEmail = authRepository.currentUserEmail
                                        viewModel.currentScreen = "menu"
                                    }
                                    .onFailure { e ->
                                        errorMessage = e.localizedMessage ?: "Błąd logowania Google"
                                    }
                                isLoading = false
                            }
                        },
                        onBiometricLogin = {
                            val activity = context.findFragmentActivity()
                            if (activity != null) {
                                biometricAuthManager.showBiometricPrompt(
                                    activity = activity,
                                    onSuccess = {
                                        scope.launch {
                                            val saved = biometricAuthManager.getSavedCredentials()
                                            if (saved != null) {
                                                isLoading = true
                                                authRepository.login(saved.first, saved.second)
                                                    .onSuccess {
                                                        currentUserEmail = authRepository.currentUserEmail
                                                        viewModel.currentScreen = "menu"
                                                    }
                                                isLoading = false
                                            }
                                        }
                                    },
                                    onError = { err -> errorMessage = err }
                                )
                            }
                        },
                        onLogout = {
                            authRepository.logout()
                            currentUserEmail = null
                            playerStats = PlayerStats()
                        },
                        onBackToMenu = {
                            errorMessage = null
                            viewModel.currentScreen = "menu"
                        }
                    )
                }
                viewModel.selectedPlaylist == null -> {
                    MainMenuScreen(viewModel = viewModel, musicPlayer = musicPlayer)
                }
                else -> {
                    GameScreen(musicPlayer = musicPlayer, viewModel = viewModel)
                }
            }
            Spacer(modifier = Modifier.height(90.dp))
        }
    }
}
\end{lstlisting}

\hspace{0.60cm}Powyższy kod (Listing~\ref{lst:main_activity}) inicjalizuje wspólną instancję odtwarzacza \texttt{MusicPlayer} oraz repozytoria chmurowe \texttt{AuthRepository} i \texttt{StatsRepository}. Efekt uboczny \texttt{LaunchedEffect(authRepository.currentUserId)} automatycznie pobiera statystyki gracza z bazy Firestore po wykryciu zmiany identyfikatora zalogowanego użytkownika. Wewnątrz kontenera \texttt{Scaffold} instrukcja warunkowa \texttt{when} steruje nawigacją pomiędzy ekranem logowania (\texttt{LoginScreen}), menu głównym (\texttt{MainMenuScreen}) oraz ekranem właściwej rozgrywki (\texttt{GameScreen}), a dolny odstęp $90\text{ dp}$ zapobiega zasłanianiu treści przez pasek nawigacyjny \texttt{BottomNavBar}.

\subsection{Implementacja repozytoriów Firebase i logowania Google (FirebaseRepositories.kt)}

\hspace{0.60cm}Warstwa dostępu do usług chmurowych została wyodrębniona do pliku \texttt{FirebaseRepositories.kt}, w którym zdefiniowano model statystyk \texttt{PlayerStats} oraz dwie wyspecjalizowane klasy repozytoriów: \texttt{AuthRepository} i \texttt{StatsRepository}.

\begin{lstlisting}[language=Kotlin, caption={Implementacja uwierzytelniania Firebase, Google Credential Manager oraz bazy Firestore w pliku FirebaseRepositories.kt}, label={lst:firebase_repositories}]
data class PlayerStats(
    val gamesPlayed: Int = 0,
    val songsGuessed: Int = 0,
    val songsMissed: Int = 0,
    val totalPoints: Int = 0,
    val bestScore: Int = 0
)

class AuthRepository(
    private val auth: FirebaseAuth = FirebaseAuth.getInstance()
) {
    val currentUserEmail: String?
        get() = auth.currentUser?.email

    val currentUserId: String?
        get() = auth.currentUser?.uid

    suspend fun login(email: String, password: String): Result<Unit> = runCatching {
        auth.signInWithEmailAndPassword(email.trim(), password).await()
    }

    suspend fun register(email: String, password: String): Result<Unit> = runCatching {
        auth.createUserWithEmailAndPassword(email.trim(), password).await()
    }

    suspend fun loginWithGoogle(context: Context): Result<Unit> = runCatching {
        val webClientIdResId = context.resources.getIdentifier("default_web_client_id", "string", context.packageName)
        val webClientId = if (webClientIdResId != 0) {
            context.getString(webClientIdResId)
        } else {
            "782578499357-fbt5i81u7b3ilf99057j23qb7t68i8p1.apps.googleusercontent.com"
        }

        val activity = context.findActivity() ?: error("Brak kontekstu Activity dla logowania Google.")
        val credentialManager = CredentialManager.create(activity)

        val googleIdOption = GetGoogleIdOption.Builder()
            .setFilterByAuthorizedAccounts(false)
            .setServerClientId(webClientId)
            .setAutoSelectEnabled(false)
            .build()

        val request = GetCredentialRequest.Builder()
            .addCredentialOption(googleIdOption)
            .build()

        val result = credentialManager.getCredential(context = activity, request = request)
        val credential = result.credential

        if (credential is androidx.credentials.CustomCredential &&
            credential.type == GoogleIdTokenCredential.TYPE_GOOGLE_ID_TOKEN_CREDENTIAL
        ) {
            val googleIdTokenCredential = GoogleIdTokenCredential.createFrom(credential.data)
            val firebaseCredential = GoogleAuthProvider.getCredential(googleIdTokenCredential.idToken, null)
            auth.signInWithCredential(firebaseCredential).await()
        } else {
            error("Nieobsługiwany typ poświadczeń Google.")
        }
    }

    fun logout() {
        auth.signOut()
    }
}

class StatsRepository(
    private val db: FirebaseFirestore = FirebaseFirestore.getInstance()
) {
    suspend fun getStats(uid: String): PlayerStats = runCatching {
        val snapshot = db.collection("users").document(uid).get().await()
        snapshot.toObject(PlayerStats::class.java) ?: PlayerStats()
    }.getOrDefault(PlayerStats())

    suspend fun recordGameResult(
        uid: String,
        guessedInSession: Int,
        missedInSession: Int,
        sessionScore: Int
    ): PlayerStats {
        val current = getStats(uid)
        val updated = PlayerStats(
            gamesPlayed = current.gamesPlayed + 1,
            songsGuessed = current.songsGuessed + guessedInSession,
            songsMissed = current.songsMissed + missedInSession,
            totalPoints = current.totalPoints + sessionScore,
            bestScore = maxOf(current.bestScore, sessionScore)
        )
        runCatching {
            db.collection("users")
                .document(uid)
                .set(updated, SetOptions.merge())
                .await()
        }
        return updated
    }
}
\end{lstlisting}

\hspace{0.60cm}Implementacja przedstawiona na Listingu~\ref{lst:firebase_repositories} wykorzystuje blok \texttt{runCatching} oraz funkcję zawieszającą \texttt{.await()} z biblioteki \texttt{kotlinx-coroutines-tasks}, co pozwala na czystą, asynchroniczną obsługę operacji sieciowych bez zagnieżdżonych wywołań zwrotnych. W metodzie \texttt{loginWithGoogle} systemowy \texttt{CredentialManager} pobiera token tożsamości Google (\texttt{GoogleIdTokenCredential}) i przekazuje go do \texttt{FirebaseAuth}. Z kolei klasa \texttt{StatsRepository} w metodzie \texttt{recordGameResult} pobiera dotychczasowy stan dokumentu gracza z kolekcji \texttt{users}, inkrementuje liczniki rozgrywek, wyznacza nowy rekord życiowy (\texttt{maxOf}) i zapisuje zaktualizowany obiekt w bazie Firestore przy użyciu opcji scalania \texttt{SetOptions.merge()}.

\subsection{Implementacja modułów sprzętowych: biometria i akcelerometr}

\hspace{0.60cm}Aplikacja wykorzystuje dwa dedykowane komponenty sprzętowe zlokalizowane w pakiecie \texttt{sensors}: moduł uwierzytelniania biometrycznego \texttt{BiometricAuthManager} oraz detektor obrotu smartfona \texttt{FlipDetector}.

\begin{lstlisting}[language=Kotlin, caption={Implementacja obsługi biometrii w pliku BiometricAuthManager.kt}, label={lst:biometric_manager}]
class BiometricAuthManager(private val context: Context) {

    private val prefs = context.getSharedPreferences("biometric_prefs", Context.MODE_PRIVATE)

    fun isBiometricAvailable(): Boolean {
        val biometricManager = BiometricManager.from(context)
        val authenticators = BiometricManager.Authenticators.BIOMETRIC_STRONG or
                BiometricManager.Authenticators.BIOMETRIC_WEAK or
                BiometricManager.Authenticators.DEVICE_CREDENTIAL
        return biometricManager.canAuthenticate(authenticators) == BiometricManager.BIOMETRIC_SUCCESS
    }

    fun saveCredentials(email: String, pass: String) {
        prefs.edit()
            .putString("bio_email", email.trim())
            .putString("bio_password", pass)
            .apply()
    }

    fun getSavedCredentials(): Pair<String, String>? {
        val email = prefs.getString("bio_email", null) ?: return null
        val pass = prefs.getString("bio_password", null) ?: return null
        return Pair(email, pass)
    }

    fun showBiometricPrompt(
        activity: FragmentActivity,
        onSuccess: () -> Unit,
        onError: (String) -> Unit
    ) {
        val executor = ContextCompat.getMainExecutor(activity)
        val biometricPrompt = BiometricPrompt(
            activity,
            executor,
            object : BiometricPrompt.AuthenticationCallback() {
                override fun onAuthenticationSucceeded(result: BiometricPrompt.AuthenticationResult) {
                    super.onAuthenticationSucceeded(result)
                    onSuccess()
                }

                override fun onAuthenticationError(errorCode: Int, errString: CharSequence) {
                    super.onAuthenticationError(errorCode, errString)
                    if (errorCode != BiometricPrompt.ERROR_USER_CANCELED &&
                        errorCode != BiometricPrompt.ERROR_NEGATIVE_BUTTON) {
                        onError("Błąd biometrii: $errString")
                    }
                }

                override fun onAuthenticationFailed() {
                    super.onAuthenticationFailed()
                    onError("Nie rozpoznano biometrii. Spróbuj ponownie.")
                }
            }
        )

        val authenticators = BiometricManager.Authenticators.BIOMETRIC_STRONG or
                BiometricManager.Authenticators.BIOMETRIC_WEAK or
                BiometricManager.Authenticators.DEVICE_CREDENTIAL

        val promptInfo = BiometricPrompt.PromptInfo.Builder()
            .setTitle("Logowanie biometryczne")
            .setSubtitle("Potwierdź tożsamość za pomocą odcisku palca, twarzy lub kodu blokady")
            .setAllowedAuthenticators(authenticators)
            .build()

        biometricPrompt.authenticate(promptInfo)
    }
}
\end{lstlisting}

\hspace{0.60cm}Klasa \texttt{BiometricAuthManager} (Listing~\ref{lst:biometric_manager}) najpierw weryfikuje za pomocą metody \texttt{isBiometricAvailable()}, czy urządzenie posiada skonfigurowany czujnik biometryczny lub blokadę ekranu. Wywołanie \texttt{showBiometricPrompt} otwiera natywne okno systemowe Androida i po pozytywnej weryfikacji tożsamości (\texttt{onAuthenticationSucceeded}) przekazuje sterowanie do funkcji zwrotnej \texttt{onSuccess}, umożliwiając automatyczne odczytanie poświadczeń z \texttt{SharedPreferences} i zalogowanie gracza do Firebase.

\begin{lstlisting}[language=Kotlin, caption={Implementacja detektora obrotu urządzenia w pliku FlipDetector.kt}, label={lst:flip_detector}]
class FlipDetector(
    sensorManager: SensorManager,
    private val onPhoneFlipped: () -> Unit
) : SensorEventListener {

    private val accelerometer: Sensor? = sensorManager.getDefaultSensor(Sensor.TYPE_ACCELEROMETER)
    private var isRegistered = false
    private var isFaceDown = false
    private var lastTriggerTime = 0L

    init {
        if (accelerometer != null) {
            isRegistered = sensorManager.registerListener(
                this,
                accelerometer,
                SensorManager.SENSOR_DELAY_UI
            )
        }
    }

    fun unregister(sensorManager: SensorManager) {
        if (isRegistered) {
            sensorManager.unregisterListener(this)
            isRegistered = false
        }
    }

    override fun onSensorChanged(event: SensorEvent?) {
        if (event?.sensor?.type != Sensor.TYPE_ACCELEROMETER) return

        val zAxis = event.values[2]
        val currentTime = System.currentTimeMillis()

        if (zAxis < -4.0f && !isFaceDown) {
            if (currentTime - lastTriggerTime > 1000L) {
                isFaceDown = true
                lastTriggerTime = currentTime
                onPhoneFlipped()
            }
        } else if (zAxis > 3.0f && isFaceDown) {
            isFaceDown = false
        }
    }

    override fun onAccuracyChanged(sensor: Sensor?, accuracy: Int) {}
}
\end{lstlisting}

\hspace{0.60cm}Moduł \texttt{FlipDetector} przedstawiony na Listingu~\ref{lst:flip_detector} rejestruje nasłuchiwanie akcelerometru z częstotliwością \texttt{SENSOR\_DELAY\_UI}. W metodzie \texttt{onSensorChanged} analizowana jest wartość składowej $Z$ wektora przyspieszenia ziemskiego: spadek poniżej wartości $-4,0\text{ m/s}^2$ przy jednoczesnym zachowaniu minimum jednosekundowego odstępu czasowego (\texttt{1000L}) wywołuje funkcję \texttt{onPhoneFlipped()}, powodując pominięcie bieżącego etapu zgadywania. Zastosowanie histerezy (odblokowanie flagi \texttt{isFaceDown} dopiero po powrocie osi $Z$ powyżej $+3,0\text{ m/s}^2$) zapobiega drganiom odczytów.

\subsection{Implementacja menu głównego i ekranu logowania}

\hspace{0.60cm}Interfejs użytkownika w warstwie prezentacji został podzielony na wyspecjalizowane ekrany kompozywalne, z których kluczowe znaczenie dla obsługi konta oraz wyboru repertuaru mają \texttt{LoginScreen.kt} oraz \texttt{MainMenuScreen.kt}.

\begin{lstlisting}[language=Kotlin, caption={Implementacja formularza logowania, rejestracji i biometrii w pliku LoginScreen.kt}, label={lst:login_screen}]
@Composable
fun LoginScreen(
    isLoading: Boolean,
    errorMessage: String?,
    currentUserEmail: String?,
    isBiometricAvailable: Boolean,
    onLogin: (String, String) -> Unit,
    onRegister: (String, String) -> Unit,
    onGoogleLogin: () -> Unit,
    onBiometricLogin: () -> Unit,
    onLogout: () -> Unit,
    onBackToMenu: () -> Unit
) {
    var email by remember { mutableStateOf("") }
    var password by remember { mutableStateOf("") }

    Column(
        modifier = Modifier
            .fillMaxSize()
            .verticalScroll(rememberScrollState())
            .padding(16.dp),
        verticalArrangement = Arrangement.Top,
        horizontalAlignment = Alignment.CenterHorizontally
    ) {
        Icon(
            imageVector = Icons.Default.AccountCircle,
            contentDescription = "Ikona logowania",
            modifier = Modifier.size(56.dp),
            tint = MaterialTheme.colorScheme.primary
        )
        Spacer(modifier = Modifier.height(12.dp))
        Text(
            text = if (currentUserEmail != null) "Zalogowano" else "Konto gracza",
            style = MaterialTheme.typography.headlineMedium,
            color = MaterialTheme.colorScheme.primary
        )
        Spacer(modifier = Modifier.height(16.dp))

        if (currentUserEmail != null) {
            Text(text = "Konto: $currentUserEmail")
            Spacer(modifier = Modifier.height(16.dp))
            Button(onClick = onLogout, enabled = !isLoading) {
                Text("Wyloguj")
            }
        } else {
            OutlinedTextField(
                value = email,
                onValueChange = { email = it },
                label = { Text("E-mail") },
                singleLine = true,
                modifier = Modifier.fillMaxWidth()
            )
            Spacer(modifier = Modifier.height(12.dp))
            OutlinedTextField(
                value = password,
                onValueChange = { password = it },
                label = { Text("Hasło (min. 6 znaków)") },
                singleLine = true,
                visualTransformation = PasswordVisualTransformation(),
                modifier = Modifier.fillMaxWidth()
            )
            Spacer(modifier = Modifier.height(16.dp))
            Button(
                onClick = { onLogin(email, password) },
                enabled = !isLoading && email.isNotBlank() && password.isNotBlank(),
                modifier = Modifier.fillMaxWidth()
            ) {
                Text("Zaloguj się")
            }
            Spacer(modifier = Modifier.height(8.dp))
            OutlinedButton(
                onClick = { onRegister(email, password) },
                enabled = !isLoading && email.isNotBlank() && password.length >= 6,
                modifier = Modifier.fillMaxWidth()
            ) {
                Text("Zarejestruj konto")
            }
            Spacer(modifier = Modifier.height(12.dp))
            HorizontalDivider(modifier = Modifier.padding(vertical = 4.dp))
            ElevatedButton(
                onClick = onGoogleLogin,
                enabled = !isLoading,
                modifier = Modifier.fillMaxWidth()
            ) {
                Icon(
                    painter = painterResource(id = R.drawable.ic_google),
                    contentDescription = "Google",
                    tint = Color.Unspecified,
                    modifier = Modifier.size(20.dp)
                )
                Spacer(modifier = Modifier.width(8.dp))
                Text("Zaloguj się przez Google")
            }

            if (isBiometricAvailable) {
                Spacer(modifier = Modifier.height(8.dp))
                OutlinedButton(
                    onClick = onBiometricLogin,
                    enabled = !isLoading,
                    modifier = Modifier.fillMaxWidth()
                ) {
                    Icon(
                        imageVector = Icons.Default.Fingerprint,
                        contentDescription = "Biometria",
                        modifier = Modifier.size(20.dp)
                    )
                    Spacer(modifier = Modifier.width(8.dp))
                    Text("Zaloguj się biometrycznie")
                }
            }
        }

        if (isLoading) {
            Spacer(modifier = Modifier.height(16.dp))
            CircularProgressIndicator()
        }
        if (errorMessage != null) {
            Spacer(modifier = Modifier.height(12.dp))
            Text(text = errorMessage, color = MaterialTheme.colorScheme.error)
        }
        Spacer(modifier = Modifier.height(24.dp))
        TextButton(onClick = onBackToMenu) {
            Text("Powrót do menu")
        }
    }
}
\end{lstlisting}

\hspace{0.60cm}Widok \texttt{LoginScreen} (Listing~\ref{lst:login_screen}) jest całkowicie bezstanowy w zakresie logiki biznesowej (wzorzec \textit{State Hoisting}) – przechowuje jedynie lokalny stan pól formularza (\texttt{email}, \texttt{password}), natomiast wszystkie zdarzenia deleguje do funkcji zwrotnych przekazanych z \texttt{MainContainer}. Pole hasła maskuje wprowadzane znaki za pomocą \texttt{PasswordVisualTransformation()}, przycisk rejestracji aktywuje się dopiero po wpisaniu minimum $6$ znaków, a przycisk logowania biometrycznego jest renderowany warunkowo wyłącznie na urządzeniach wspierających tę technologię (\texttt{isBiometricAvailable}).

\begin{lstlisting}[language=Kotlin, caption={Implementacja wyszukiwarki i listy playlist w pliku MainMenuScreen.kt}, label={lst:main_menu_screen}]
@Composable
fun MainMenuScreen(
    viewModel: GameViewModel,
    musicPlayer: MusicPlayer
) {
    Column(
        modifier = Modifier
            .fillMaxSize()
            .padding(24.dp),
        verticalArrangement = Arrangement.Top,
        horizontalAlignment = Alignment.CenterHorizontally
    ) {
        OutlinedTextField(
            value = viewModel.searchQuery,
            onValueChange = { viewModel.updateSearchQuery(it) },
            label = { Text("Wyszukaj playliste (np. Rock)") },
            singleLine = true,
            modifier = Modifier.fillMaxWidth(),
            shape = RoundedCornerShape(16.dp),
            trailingIcon = {
                if (viewModel.searchQuery.isNotEmpty()) {
                    IconButton(onClick = { viewModel.clearSearch() }) {
                        Icon(Icons.Default.Close, contentDescription = "Wyczysc")
                    }
                }
            }
        )

        Spacer(modifier = Modifier.height(8.dp))

        Button(
            onClick = { viewModel.searchForPlaylists(musicPlayer) },
            modifier = Modifier.fillMaxWidth(),
            enabled = viewModel.searchQuery.isNotBlank() && !viewModel.isSearching
        ) {
            if (viewModel.isSearching) {
                CircularProgressIndicator(modifier = Modifier.size(24.dp))
            } else {
                Icon(Icons.Default.Search, contentDescription = null)
                Spacer(modifier = Modifier.width(8.dp))
                Text("Szukaj w Deezer")
            }
        }

        val listToShow = if (viewModel.searchResults.isNotEmpty()) {
            viewModel.searchResults
        } else {
            viewModel.availablePlaylists
        }

        Column(modifier = Modifier.fillMaxWidth()) {
            listToShow.forEach { playlist ->
                Card(
                    modifier = Modifier
                        .fillMaxWidth()
                        .padding(vertical = 4.dp)
                        .clickable { viewModel.selectPlaylist(playlist) },
                    shape = RoundedCornerShape(16.dp)
                ) {
                    Column(modifier = Modifier.padding(16.dp)) {
                        Text(text = playlist.name, style = MaterialTheme.typography.titleMedium)
                        Text(text = playlist.description, style = MaterialTheme.typography.bodyMedium)
                    }
                }
            }
        }
    }
}
\end{lstlisting}

\hspace{0.60cm}Implementacja przedstawiona na Listingu~\ref{lst:main_menu_screen} wykorzystuje pole tekstowe \texttt{OutlinedTextField} z dynamiczną ikoną czyszczenia wpisanej frazy (\texttt{trailingIcon}). Poniżej komponentu wyszukiwarki system sprawdza zawartość listy \texttt{searchResults} – jeżeli użytkownik nie wykonał jeszcze własnego zapytania, wyświetlane są domyślne kategorie zdefiniowane w \texttt{availablePlaylists}, a kliknięcie dowolnej karty (\texttt{Card}) natychmiast inicjuje rozgrywkę dla wybranej listy.

\subsection{Implementacja warstwy sieciowej i audio (DeezerAPI.kt)}

\hspace{0.60cm}Poniższy listing przedstawia implementację klasy \texttt{MusicPlayer}, w szczególności asynchroniczne ładowanie i odtwarzanie próbki utworu z wykorzystaniem \texttt{prepareAsync()} oraz obsługę limitu czasu odtwarzania za pomocą obiektu \texttt{Handler}:

\begin{lstlisting}[language=Kotlin, caption={Implementacja odtwarzania próbek i komunikacji z Deezer API w pliku DeezerAPI.kt}, label={lst:deezer_api_real}]
data class SongSuggestion(
    val id: Int,
    val title: String,
    val artist: String
)

class MusicPlayer {
    private var mediaPlayer: MediaPlayer? = null
    private val handler = Handler(Looper.getMainLooper())

    suspend fun loadAndPlaySong(
        trackId: Int,
        durationMillis: Long,
        onSongFinished: () -> Unit,
        onNetworkError: () -> Unit
    ): Triple<String, String, String>? {
        stop()
        return try {
            val jsonString = withContext(Dispatchers.IO) {
                URL("https://api.deezer.com/track/$trackId?ts=${System.currentTimeMillis()}").readText()
            }
            val jsonObject = JSONObject(jsonString)
            val previewUrl = jsonObject.optString("preview", "")

            if (previewUrl.isEmpty()) return null

            val title = jsonObject.getString("title")
            val artist = jsonObject.getJSONObject("artist").getString("name")
            val coverUrl = jsonObject.getJSONObject("album").optString("cover_medium", "")

            val newPlayer = MediaPlayer().apply {
                setDataSource(previewUrl)
                setOnPreparedListener { player ->
                    player.start()
                    handler.postDelayed({
                        if (isPlaying) {
                            pause()
                            onSongFinished()
                        }
                    }, durationMillis)
                }
                setOnErrorListener { _, _, _ ->
                    onNetworkError()
                    true
                }
                prepareAsync()
            }
            mediaPlayer = newPlayer
            Triple(title, artist, coverUrl)
        } catch (e: Exception) {
            onNetworkError()
            null
        }
    }

    fun playSnippet(durationMillis: Long) {
        mediaPlayer?.let { player ->
            handler.removeCallbacksAndMessages(null)
            player.seekTo(0)
            player.start()
            handler.postDelayed({
                if (player.isPlaying) player.pause()
            }, durationMillis)
        }
    }
}
\end{lstlisting}

\hspace{0.60cm}Metoda \texttt{loadAndPlaySong} zaprezentowana na Listingu~\ref{lst:deezer_api_real} wykonuje zapytanie sieciowe w wątku roboczym \texttt{Dispatchers.IO}, co chroni główny wątek aplikacji przed zawieszeniem. Po wyodrębnieniu adresu pliku MP3 z obiektu JSON odtwarzacz jest buforowany asynchronicznie (\texttt{prepareAsync()}), a w momencie rozpoczęcia dźwięku obiekt \texttt{Handler} planuje opóźnione wywołanie funkcji pauzującej (\texttt{postDelayed}) dokładnie po upływie czasu zadanego dla bieżącego etapu. Metoda \texttt{playSnippet} umożliwia ponowne odsłuchanie tego samego fragmentu od początku bez konieczności ponownego pobierania pliku z sieci.

\subsection{Implementacja logiki rozgrywki i punktacji (GameViewModel.kt)}

\hspace{0.60cm}W klasie \texttt{GameViewModel} zdefiniowano zarządzanie pięciorundową sesją gry, system punktacji malejącej zależnej od numeru próby (\texttt{attemptScores}), obsługę pomijania pojedynczego etapu (\texttt{skipCurrentAttempt}) oraz mechanizm bezpiecznego opuszczania trwającej rozgrywki (\texttt{requestExit}).

\begin{lstlisting}[language=Kotlin, caption={Kluczowe metody logiki rozgrywki, punktacji i nawigacji w klasie GameViewModel.kt}, label={lst:game_viewmodel_real}]
class GameViewModel : ViewModel() {
    var score by mutableStateOf(0)
        private set
    var currentRound by mutableStateOf(1)
        private set
    val totalRounds = 5
    var attemptsLeft by mutableStateOf(6)
        private set
    val incorrectGuesses = mutableStateListOf<String>()

    private val attemptDurations = listOf(200L, 1000L, 3000L, 6000L, 10000L, 15000L)
    private val attemptScores = listOf(5000, 4000, 3000, 2000, 1000, 500)
    val maxScore = 25000

    var showExitConfirmation by mutableStateOf(false)
        private set
    private var pendingDestination: String = "menu"
    var currentScreen by mutableStateOf("menu")

    fun getCurrentSnippetDuration(): Long {
        val attemptIndex = 6 - attemptsLeft
        return if (attemptIndex in attemptDurations.indices) {
            attemptDurations[attemptIndex]
        } else {
            15000L
        }
    }

    fun checkAnswer(musicPlayer: MusicPlayer) {
        if (isRoundOver || attemptsLeft <= 0 || isLoadingSong) return
        val currentGuess = userGuess.trim()
        if (currentGuess.isEmpty()) return

        songSuggestions.clear()
        showSuggestions = false

        if (currentGuess.equals(correctTitle.trim(), ignoreCase = true)) {
            val attemptIndex = 6 - attemptsLeft
            val earnedPoints = if (attemptIndex in attemptScores.indices) attemptScores[attemptIndex] else 0

            score += earnedPoints
            isAnswerCorrect = true
            isRoundOver = true
            isPlaying = false
            musicPlayer.stop()
            message = "Brawo! Poprawna odpowiedź (+${earnedPoints} pkt)."
        } else {
            attemptsLeft -= 1
            incorrectGuesses.add(currentGuess)
            if (incorrectGuesses.size > 5) incorrectGuesses.removeAt(0)
            userGuess = ""

            if (attemptsLeft > 0) {
                val nextDuration = getCurrentSnippetDuration()
                musicPlayer.playSnippet(nextDuration)
            } else {
                isAnswerCorrect = false
                isRoundOver = true
                isPlaying = false
                musicPlayer.stop()
                message = "Koniec prób! 0 punktów. Prawidłowy tytuł to: $correctTitle"
            }
        }
    }

    fun skipCurrentAttempt(musicPlayer: MusicPlayer) {
        if (isRoundOver || isLoadingSong) return
        if (attemptsLeft > 1) {
            attemptsLeft -= 1
            val nextDuration = getCurrentSnippetDuration()
            musicPlayer.playSnippet(nextDuration)
            isPlaying = true
        } else {
            skipSong(musicPlayer)
        }
    }

    fun requestExit(destination: String) {
        if (selectedPlaylist != null && !isGameSessionFinished) {
            pendingDestination = destination
            showExitConfirmation = true
        } else {
            currentScreen = destination
        }
    }

    fun confirmExit(musicPlayer: MusicPlayer) {
        showExitConfirmation = false
        currentScreen = pendingDestination
        if (pendingDestination == "menu") {
            resetPlaylistSelection(musicPlayer)
        }
    }
}
\end{lstlisting}

\hspace{0.60cm}Kod na Listingu~\ref{lst:game_viewmodel_real} implementuje system nagradzania gracza za szybkie rozpoznanie utworu: na podstawie indeksu bieżącej próby (\texttt{6 - attemptsLeft}) metoda \texttt{checkAnswer} pobiera z tablicy \texttt{attemptScores} odpowiednią liczbę punktów (od $5000$ pkt w pierwszej próbie do $500$ pkt w szóstej). Metoda \texttt{skipCurrentAttempt} odpowiada zarówno za przycisk pominięcia etapu na ekranie, jak i za gest odwrócenia telefonu wykrywany przez \texttt{FlipDetector}. Dodatkowo metoda \texttt{requestExit} chroni gracza przed przypadkową utratą postępu – jeżeli użytkownik kliknie przycisk na dolnym pasku nawigacyjnym w trakcie aktywnej rozgrywki, system wstrzymuje przejście i wyświetla okno dialogowe potwierdzenia (\texttt{showExitConfirmation = true}).

\subsection{Implementacja interfejsu rozgrywki (GameScreen.kt)}

\hspace{0.60cm}Ekran rozgrywki w pliku \texttt{GameScreen.kt} dynamicznie dostosowuje wyświetlane komponenty do aktualnego etapu rundy. Istotnym elementem implementacji jest ograniczenie maksymalnej wysokości karty z podpowiedziami wyszukiwarki za pomocą modyfikatora \texttt{.heightIn(max = 180.dp)} połączonego z pionowym przewijaniem (\texttt{verticalScroll}), co zapobiega rozpychaniu ekranu po wysunięciu klawiatury systemowej.

\begin{lstlisting}[language=Kotlin, caption={Fragment implementacji listy podpowiedzi oraz paska postępu w pliku GameScreen.kt}, label={lst:game_screen_real}]
val currentAttemptIndex = 6 - viewModel.attemptsLeft
val progress = (currentAttemptIndex + 1).toFloat() / 6f

LinearProgressIndicator(
    progress = { progress },
    modifier = Modifier
        .fillMaxWidth()
        .height(8.dp)
        .clip(RoundedCornerShape(4.dp)),
    color = MaterialTheme.colorScheme.primary,
    trackColor = MaterialTheme.colorScheme.surfaceVariant,
)

if (viewModel.showSuggestions && viewModel.songSuggestions.isNotEmpty()) {
    Card(
        modifier = Modifier.fillMaxWidth().padding(bottom = 8.dp),
        shape = RoundedCornerShape(12.dp)
    ) {
        Column(
            modifier = Modifier
                .fillMaxWidth()
                .heightIn(max = 180.dp)
                .verticalScroll(rememberScrollState())
                .padding(4.dp)
        ) {
            viewModel.songSuggestions.forEach { suggestion ->
                Column(
                    modifier = Modifier
                        .fillMaxWidth()
                        .clickable { viewModel.selectSuggestion(suggestion, musicPlayer) }
                        .padding(10.dp)
                ) {
                    Text(text = suggestion.title, fontWeight = FontWeight.Bold)
                    Text(text = suggestion.artist, fontSize = 12.sp)
                }
            }
        }
    }
}
\end{lstlisting}

\hspace{0.60cm}Fragment przedstawiony na Listingu~\ref{lst:game_screen_real} wizualizuje postęp zużytych prób za pomocą paska \texttt{LinearProgressIndicator}, którego wypełnienie rośnie proporcjonalnie od $1/6$ do $6/6$ wraz z kolejnymi etapami odsłuchu. Poniżej renderowana jest karta podpowiedzi pobieranych na żywo z API Deezer – zastosowanie przewijanej kolumny o maksymalnej wysokości $180\text{ dp}$ sprawia, że lista sugestii pozostaje czytelna i dostępna do kliknięcia nawet przy otwartej klawiaturze ekranowej smartfona.

\subsection{Przebieg zmian i rozwoju aplikacji w ostatnich tygodniach}

\hspace{0.60cm}W trakcie prac nad projektem architektura aplikacji przeszła szereg iteracyjnych ulepszeń:
\begin{enumerate}
    \item \textbf{Integracja z usługami Firebase (Auth oraz Cloud Firestore):} Dodano pełną obsługę kont użytkowników (rejestrację i logowanie za pomocą adresu e-mail oraz hasła) oraz trwały zapis statystyk rozgrywki i rekordów punktowych w bazie dokumentowej Firestore.
    \item \textbf{Wdrożenie logowania Google oraz uwierzytelniania biometrycznego:} Zintegrowano systemowy mechanizm \texttt{CredentialManager} do logowania kontem Google oraz bibliotekę \texttt{BiometricPrompt}, umożliwiającą szybkie logowanie odciskiem palca lub skanem twarzy.
    \item \textbf{Rozbudowa katalogu muzycznego o listy przebojów (\textit{Charts}) i system punktacji malejącej:} Rozszerzono integrację z Deezer API o automatyczne pobieranie topowych playlist, artystów i albumów, a także wdrożono gradację punktów ($5000$–$500$ pkt) premiującą jak najszybsze rozpoznanie utworu.
    \item \textbf{Uporządkowanie struktury pakietów i nawigacji:} Podzielono kod na warstwy funkcjonalne (\texttt{data/firebase}, \texttt{data/remote}, \texttt{sensors}, \texttt{ui/screens}, \texttt{ui/components}, \texttt{viewmodel}), wdrożono ciemny motyw z neonowymi akcentami (\texttt{NeonBlue}) oraz dodano okno potwierdzenia wyjścia chroniące przed przypadkowym przerwaniem gry.
\end{enumerate}